% 第3章
\part{スクレーパ機構の統合}
\label{scraper}

\section{剥離ユニットの構成}
\subsection{リニアアクチュエータによる昇降・押し込み機構}
剥離機構には，
リニアアクチュエータを用いたシリンダ駆動式スクレーパを採用した．
これは，無線コントローラからの信号により，
刃先の昇降および壁面への押し込み動作を任意に行えるものである．
作業地点までは刃を上げ，剥離開始時に強力に塗膜へ貫入させる制御を可能とした．

\begin{figure}[H] 
  \centering 
  \includegraphics[width=0.6\textwidth]{tougou.png} 
  \vspace{-2mm} 
  \caption{Scraper unit integrated with DC linear actuator} \label{fig:result} 
  \end{figure}
\newpage

\section{刃先角度の設定と構造的制約}
\subsection{角度40度の採用理由}
剥離性能において，
刃先が壁面となす角度は極めて重要である．
本機体では，現行の機体フレームへの統合構造とリニアアクチュエータの可動域の兼ね合いから，
壁面接地時の刃角が約40度に固定される設計となっている．

\subsection{40度設定の力学的影響}
一般に，40度という深い角度は塗膜への「初期貫入（刺さる動作）」には非常に有利である．
しかし，貫入した状態で水平に走行しようとすると，
壁面から機体を押し返そうとする垂直方向の力（法線抗力）が大きくなる傾向がある．
この構造的制約が実際の走行にどのような影響を与えたかについて，次章で詳述する．

\newpage